\centerline{\textbf{\large Dinámica de los Precios en el Sector de Materiales de }}
\centerline{\textbf{\large Construcción: Patrones y Predicciones Basados en Datos}}

\vspace{2cm}

\rightline{Autor: Diego Milciades Espínola Alvarenga}
\rightline{Matías Rubén Sánchez Sánchez}
\vspace{0.5cm}
\rightline{Asesor: D.Sc. Gustavo Sosa-Cabrera}

\vspace{0.5cm}

\noindent
\textbf{RESUMEN}

\vspace{0.5cm}

\noindent
\justify
El sector de la construcción en Paraguay enfrenta un problema concreto: los precios de sus insumos cambian de manera imprevisible y los actores del mercado no cuentan con herramientas para anticiparlo. Esta investigación aborda esa brecha centrándose en dos materiales de amplio consumo —cemento Yguazú y ladrillo común Tobatí— a partir de registros mensuales del período 2014--2025. Se incorporaron además dos variables exógenas de relevancia local: el confinamiento por COVID-19 y el nivel del río Paraguay en el puerto de Asunción, cuya bajante histórica de 2021 obligó a la Industria Nacional del Cemento a sustituir el transporte fluvial por vía terrestre, encareciendo los costos de distribución.
\justify
El proceso de modelado partió de la recolección y limpieza de las series históricas, incluyendo el tratamiento de valores faltantes mediante interpolación polinómica cuadrática, seleccionada por ajustarse mejor a los valores de referencia disponibles. Sobre esa base se construyeron tres enfoques predictivos: el modelo estadístico SARIMAX y dos arquitecturas de redes neuronales recurrentes, LSTM y GRU. Todos se evaluaron con un horizonte de 24 meses bajo escenarios alternativos de variables exógenas, usando el RMSE como métrica central de comparación.
\justify
Los modelos neuronales superaron al enfoque estadístico. El LSTM multivariado registró el mejor desempeño global, mientras que el GRU alcanzó resultados comparables con menor costo computacional. SARIMAX, aunque útil como línea de base, no logró capturar la no linealidad característica de estas series. Un hallazgo relevante fue que predecir previamente el nivel del río con un LSTM univariado y luego usar esa proyección como variable exógena mejoró la precisión de los modelos de precios, validando la hipótesis de que el factor hidrológico incide sobre la dinámica de costos del sector.
\justify
El resultado final es un sistema de predicción replicable, construido íntegramente con datos públicos y herramientas de acceso libre. Su aplicación puede extenderse a otros materiales o indicadores económicos del sector construcción en Paraguay.

\vspace{0.5cm}

\noindent
\textbf{Palabras Clave}: Precios de materiales de construcción,variables exógenas, río Paraguay, COVID-19, modelos predictivos, SARIMAX, LSTM, GRU, Paraguay.
